% 附录E：坐标系与微分算子

\chapter{坐标系与微分算子}

\section{柱坐标系 $(r, \varphi, z)$}

\textbf{坐标关系}：
\begin{align}
    x &= r\cos\varphi, \quad y = r\sin\varphi, \quad z = z
\end{align}

\textbf{逆关系}：
\begin{align}
    r &= \sqrt{x^2 + y^2}, \quad \varphi = \arctan\frac{y}{x}, \quad z = z
\end{align}

\textbf{体积元}：
\begin{equation}
    \mathrm{d}V = r\,\mathrm{d}r\,\mathrm{d}\varphi\,\mathrm{d}z
\end{equation}

\textbf{面积元}：
\begin{align}
    \mathrm{d}S_r &= r\,\mathrm{d}\varphi\,\mathrm{d}z \quad \text{（$r=$常数面上）} \\
    \mathrm{d}S_z &= r\,\mathrm{d}r\,\mathrm{d}\varphi \quad \text{（$z=$常数面上）}
\end{align}

\textbf{线元}：
\begin{equation}
    \mathrm{d}l^2 = \mathrm{d}r^2 + r^2\mathrm{d}\varphi^2 + \mathrm{d}z^2
\end{equation}

\textbf{拉普拉斯算子}：
\begin{zhongyao}[柱坐标系拉普拉斯算子]
\[
    \nabla^2 u = \frac{1}{r}\frac{\partial}{\partial r}\left(r\frac{\partial u}{\partial r}\right) + \frac{1}{r^2}\frac{\partial^2 u}{\partial\varphi^2} + \frac{\partial^2 u}{\partial z^2}
\]
\end{zhongyao}

\section{球坐标系 $(r, \theta, \varphi)$}

\textbf{坐标关系}：
\begin{align}
    x &= r\sin\theta\cos\varphi \\
    y &= r\sin\theta\sin\varphi \\
    z &= r\cos\theta
\end{align}

\textbf{逆关系}：
\begin{align}
    r &= \sqrt{x^2 + y^2 + z^2} \\
    \theta &= \arccos\frac{z}{\sqrt{x^2+y^2+z^2}} \\
    \varphi &= \arctan\frac{y}{x}
\end{align}

\textbf{体积元}：
\begin{equation}
    \mathrm{d}V = r^2\sin\theta\,\mathrm{d}r\,\mathrm{d}\theta\,\mathrm{d}\varphi
\end{equation}

\textbf{面积元}：
\begin{align}
    \mathrm{d}S_r &= r^2\sin\theta\,\mathrm{d}\theta\,\mathrm{d}\varphi \quad \text{（$r=$常数面上）} \\
    \mathrm{d}S_\theta &= r\sin\theta\,\mathrm{d}r\,\mathrm{d}\varphi \quad \text{（$\theta=$常数面上）} \\
    \mathrm{d}S_\varphi &= r\,\mathrm{d}r\,\mathrm{d}\theta \quad \text{（$\varphi=$常数面上）}
\end{align}

\textbf{线元}：
\begin{equation}
    \mathrm{d}l^2 = \mathrm{d}r^2 + r^2\mathrm{d}\theta^2 + r^2\sin^2\theta\,\mathrm{d}\varphi^2
\end{equation}

\textbf{拉普拉斯算子}：
\begin{zhongyao}[球坐标系拉普拉斯算子]
\[
    \nabla^2 u = \frac{1}{r^2}\frac{\partial}{\partial r}\left(r^2\frac{\partial u}{\partial r}\right) + \frac{1}{r^2\sin\theta}\frac{\partial}{\partial\theta}\left(\sin\theta\frac{\partial u}{\partial\theta}\right) + \frac{1}{r^2\sin^2\theta}\frac{\partial^2 u}{\partial\varphi^2}
\]
\end{zhongyao}

\section{梯度、散度、旋度}

\subsection{直角坐标系}
\begin{align}
    \nabla f &= \frac{\partial f}{\partial x}\boldsymbol{i} + \frac{\partial f}{\partial y}\boldsymbol{j} + \frac{\partial f}{\partial z}\boldsymbol{k} \\
    \nabla \cdot \boldsymbol{F} &= \frac{\partial F_x}{\partial x} + \frac{\partial F_y}{\partial y} + \frac{\partial F_z}{\partial z} \\
    \nabla \times \boldsymbol{F} &= \left(\frac{\partial F_z}{\partial y} - \frac{\partial F_y}{\partial z}\right)\boldsymbol{i} + \left(\frac{\partial F_x}{\partial z} - \frac{\partial F_z}{\partial x}\right)\boldsymbol{j} + \left(\frac{\partial F_y}{\partial x} - \frac{\partial F_x}{\partial y}\right)\boldsymbol{k}
\end{align}

\subsection{柱坐标系}
\begin{align}
    \nabla f &= \frac{\partial f}{\partial r}\boldsymbol{e}_r + \frac{1}{r}\frac{\partial f}{\partial\varphi}\boldsymbol{e}_\varphi + \frac{\partial f}{\partial z}\boldsymbol{e}_z \\
    \nabla \cdot \boldsymbol{F} &= \frac{1}{r}\frac{\partial}{\partial r}(rF_r) + \frac{1}{r}\frac{\partial F_\varphi}{\partial\varphi} + \frac{\partial F_z}{\partial z} \\
    \nabla \times \boldsymbol{F} &= \left(\frac{1}{r}\frac{\partial F_z}{\partial\varphi} - \frac{\partial F_\varphi}{\partial z}\right)\boldsymbol{e}_r + \left(\frac{\partial F_r}{\partial z} - \frac{\partial F_z}{\partial r}\right)\boldsymbol{e}_\varphi + \frac{1}{r}\left[\frac{\partial}{\partial r}(rF_\varphi) - \frac{\partial F_r}{\partial\varphi}\right]\boldsymbol{e}_z
\end{align}

\subsection{球坐标系}
\begin{align}
    \nabla f &= \frac{\partial f}{\partial r}\boldsymbol{e}_r + \frac{1}{r}\frac{\partial f}{\partial\theta}\boldsymbol{e}_\theta + \frac{1}{r\sin\theta}\frac{\partial f}{\partial\varphi}\boldsymbol{e}_\varphi \\
    \nabla \cdot \boldsymbol{F} &= \frac{1}{r^2}\frac{\partial}{\partial r}(r^2F_r) + \frac{1}{r\sin\theta}\frac{\partial}{\partial\theta}(\sin\theta F_\theta) + \frac{1}{r\sin\theta}\frac{\partial F_\varphi}{\partial\varphi} \\
    \nabla \times \boldsymbol{F} &= \frac{1}{r\sin\theta}\left[\frac{\partial}{\partial\theta}(\sin\theta F_\varphi) - \frac{\partial F_\theta}{\partial\varphi}\right]\boldsymbol{e}_r \\
    &\quad + \frac{1}{r}\left[\frac{1}{\sin\theta}\frac{\partial F_r}{\partial\varphi} - \frac{\partial}{\partial r}(rF_\varphi)\right]\boldsymbol{e}_\theta \\
    &\quad + \frac{1}{r}\left[\frac{\partial}{\partial r}(rF_\theta) - \frac{\partial F_r}{\partial\theta}\right]\boldsymbol{e}_\varphi
\end{align}

\section{矢量恒等式}
\begin{align}
    \nabla \times (\nabla f) &= 0 \\
    \nabla \cdot (\nabla \times \boldsymbol{F}) &= 0 \\
    \nabla \times (\nabla \times \boldsymbol{F}) &= \nabla(\nabla \cdot \boldsymbol{F}) - \nabla^2 \boldsymbol{F} \\
    \nabla \cdot (f\boldsymbol{F}) &= f\nabla \cdot \boldsymbol{F} + \boldsymbol{F} \cdot \nabla f \\
    \nabla \times (f\boldsymbol{F}) &= f\nabla \times \boldsymbol{F} + \nabla f \times \boldsymbol{F}
\end{align}
